\documentclass{beamer}
\usepackage{amsmath}
\setbeameroption{show notes on second screen=right}
\setbeamertemplate{note page}[plain]

\beamertemplatenavigationsymbolsempty

\begin{document}

\begin{frame}
\frametitle{}

To divide two fractions, we multiply the first fraction by the reciprocal of the second fraction.
\note[item]{\alert<1>{To divide two fractions, we multiply the first fraction by the reciprocal of the second fraction.}}
\pause Recall when multiplying fractions, there is no requirement for a common denominator.
\note[item]{\alert<2>{Recall when multiplying fractions, there is no requirement for a common denominator.}}
\note[item]{\alert<3>{To multiply fractions, we just multiply straight across.}}
\note[item]{\alert<4>{So, we are able to convert a problem about dividing fractions into a problem about multiplying fractions.}}

\vskip10pt

\pause 
\pause 
\pause 
{\bf Example.} $\frac25 \div \frac37$.
\note[item]{\alert<5>{Let's look at an example. What is two fifths divided by three sevenths?}}

\vskip4pt

\pause $=\frac{2}{5} \times \frac{7}{3}$
\note[item]{\alert<6>{Replace dividing by three sevenths instead with multiplying by seven thirds.}}

\vskip4pt

\pause $=\frac{2 \cdot 7}{5 \cdot 3}$ 
\note[item]{\alert<7>{Now that we are multiplying fractions, we just multiply straight across. Two times seven on top, and five times three on bottom.}}

\vskip4pt

\pause $=\frac{14}{15}$
\note[item]{\alert<8>{Simplifying the top and bottom, our single fraction is fourteen over fifteen.}}

\end{frame}

\begin{frame}
\frametitle{}

{\bf Example.} $\dfrac{\phantom{.}\frac{3}{10}\phantom{.}}{\frac{7}{11}}$. 
\note[item]{\alert<1>{Let's look at another example. What is three tenths divided by seven elevenths?}}
\note[item]{\alert<2>{The format here looks different, but this is the same as if we wrote three tenths and a traditional division symbol and then seven elevenths.}}

\vskip4pt

\pause
\pause $=\frac{3}{10} \times \frac{11}{7}$
\note[item]{\alert<3>{Replace dividing by seven elevenths instead with multiplying by eleven sevenths.}}


\vskip4pt

\pause $=\frac{3 \cdot 11}{10 \cdot 7}$
\note[item]{\alert<4>{Now that we are multiplying fractions, multiply straight across. Remember, there is no need for a common denominator. So, 3 times 11 on top, and 10 times 7 on bottom.}}

\vskip4pt

\pause $=\frac{33}{70}$
\note[item]{\alert<5>{Simplifying the top and bottom, we have 33 over 70.}}


\end{frame}



%% Blank frames at the end to prevent weird shifts.
\begin{frame}
\end{frame}
\begin{frame}
\end{frame}
\begin{frame}
\end{frame}
\begin{frame}
\end{frame}
\begin{frame}
\end{frame}
\begin{frame}
\end{frame}
\begin{frame}
\end{frame}

\end{document}
